%% ma: L12-B04-0009
%% lop: 12
%% chuong: 1
%% bai: 4
%% dang: 12.4.9
%% loai: TN4
%% muc: TH
%% dap_an: "B"
%% nguon: "(NBV)-ĐỀ SỐ 11-MỖI NGÀY 1 ĐỀ THI 2025.pdf (D11, MỖI NGÀY 1 ĐỀ - Đề số 11), Phần I Câu 5"
%% kien_thuc: "Bài 4 (tâm đối xứng của đồ thị hàm phân thức bậc hai trên bậc nhất), Bài 3 (tiệm cận đứng, tiệm cận xiên)"
%% trang_thai: chinh-thuc
%% ghi_chu: "Đã duyệt 13/10/2026. Đáp án gốc B đúng."
\begin{ex}
Đồ thị của hàm số $y=\dfrac{-x^2-x+2}{x+1}$ có đặc điểm nào sau đây?
\choice{Nhận điểm $(1;1)$ làm tâm đối xứng}{\True Nhận điểm $(-1;1)$ làm tâm đối xứng}{Nhận điểm $(-1;-1)$ làm tâm đối xứng}{Không có tâm đối xứng}
\loigiai{Ta có $y=-x+\dfrac{2}{x+1}$. Đồ thị có tiệm cận đứng $x=-1$ và tiệm cận xiên $y=-x$. Giao điểm của hai tiệm cận là $I(-1;1)$ nên $I$ là tâm đối xứng của đồ thị. Chọn B.}
\end{ex}
